% corpus_supplement: Erdős problem #699
% source: https://www.erdosproblems.com/latex/699
% fetched_at: 2026-07-14T18:48:21Z
% payload_sha256: 1c0e5671725e9e173cfb17806c584973949adabd225addc507980793b6e04202
% statement/remarks/references extracted by erdos_ingest.extract_source_page;
% rendered by erdos_ingest.render_tex — do not edit
\documentclass{article}
\usepackage{amsmath, amssymb, amsthm}
\usepackage{hyperref}
\usepackage{geometry}
\geometry{margin=1in}

\title{Erd\H{o}s Problem \#699}
\date{Last updated: 2025-08-31}
\author{Source: erdosproblems.com}

\begin{document}
\maketitle

\noindent\textbf{Status:} OPEN \quad \textbf{No prize}

\noindent\textbf{Tags:} number theory, binomial coefficients

\medskip
\noindent\textbf{Problem Statement:}

Is it true that for every $1\leq i<j\leq n/2$ there exists some prime $p\geq i$ such that\[p\mid \textrm{gcd}\left(\binom{n}{i}, \binom{n}{j}\right)?\]

\bigskip
\noindent\textbf{Remarks:}

A problem of Erd\H{o}s and Szekeres. A theorem of Sylvester and Schur says that for any $1\leq i\leq n/2$ there exists some prime $p>i$ which divides $\binom{n}{i}$.

Erd\H{o}s and Szekeres further conjectured that $p\geq i$ can be improved to $p>i$ except in a few special cases. In particular this fails when $i=2$ and $n$ being some particular powers of $2$. They also found some counterexamples when $i=3$, but only one counterexample when $i\geq 4$:\[\textrm{gcd}\left(\binom{28}{5},\binom{28}{14}\right)=2^3\cdot 3^3\cdot 5.\]This is mentioned in problem B31 of Guy's collection \cite{Gu04}.

\bigskip
\noindent\textbf{References:}

[Gu04] Guy, Richard K., Unsolved problems in number theory. (2004), xviii+437.

\bigskip
\noindent\small{Source: \url{https://www.erdosproblems.com/699}}
\end{document}
